\documentclass[border=5pt]{standalone}
\usepackage{tikz}
\usetikzlibrary{arrows.meta, calc, decorations.markings}

\begin{document}

\definecolor{axisX}{RGB}{220, 50, 30}   % Red - X forward (North)
\definecolor{axisY}{RGB}{40, 160, 40}   % Green - Y right (East)
\definecolor{axisZ}{RGB}{0, 120, 200}   % Blue - Z down
\definecolor{bodycolor}{RGB}{80, 80, 80}

\begin{tikzpicture}[
    axis/.style={-{Stealth[length=5mm, width=3mm]}, line width=2.5pt},
    rotarrow/.style={
      -{Stealth[length=3.5mm, width=2.5mm]},
      line width=1.8pt,
    },
  ]

  % Isometric unit vectors
  \pgfmathsetmacro{\Xx}{-0.71}
  \pgfmathsetmacro{\Xy}{0.50}
  \pgfmathsetmacro{\Yx}{1.0}
  \pgfmathsetmacro{\Yy}{0.0}
  \pgfmathsetmacro{\Zx}{0.0}
  \pgfmathsetmacro{\Zy}{-1.0}

  % Body - isometric parallelogram
  \fill[bodycolor!10]
    ({1.1*\Xx+1.1*\Yx},{1.1*\Xy+1.1*\Yy}) --
    ({1.1*\Xx-1.1*\Yx},{1.1*\Xy-1.1*\Yy}) --
    ({-1.1*\Xx-1.1*\Yx},{-1.1*\Xy-1.1*\Yy}) --
    ({-1.1*\Xx+1.1*\Yx},{-1.1*\Xy+1.1*\Yy}) -- cycle;
  \draw[bodycolor, line width=1.5pt]
    ({1.1*\Xx+1.1*\Yx},{1.1*\Xy+1.1*\Yy}) --
    ({1.1*\Xx-1.1*\Yx},{1.1*\Xy-1.1*\Yy}) --
    ({-1.1*\Xx-1.1*\Yx},{-1.1*\Xy-1.1*\Yy}) --
    ({-1.1*\Xx+1.1*\Yx},{-1.1*\Xy+1.1*\Yy}) -- cycle;

  % Axis length
  \pgfmathsetmacro{\AL}{4.5}

  % X axis (Forward / North)
  \draw[axis, axisX] (0,0) -- ({\AL*\Xx},{\AL*\Xy});

  % Front marker triangle -- drawn AFTER the X axis line so its opaque fill
  % paints over the line where it would otherwise cross the marker,
  % keeping the marker readable as one solid shape (its base midpoint sits
  % exactly on the X-axis line, which used to bisect it when drawn first).
  % 前方マーカーの三角形。X軸の線より後に描く（三角形の塗りが不透明なため、
  % そのままだと線が三角形を貫通する箇所を上塗りし、マーカーが1つの
  % 塗りつぶし形状として読めるようにする）。三角形の底辺の中点はちょうど
  % X軸の直線上にあり、先に描くと三角形が線で分断されて見えていた。
  \fill[axisX!50]
    ({1.1*\Xx-0.35},{1.1*\Xy}) --
    ({1.1*\Xx+0.35},{1.1*\Xy}) --
    ({1.1*\Xx},{1.1*\Xy+0.3}) -- cycle;

  \node[font=\Large\bfseries, axisX, above left]
    at ({\AL*\Xx},{\AL*\Xy}) {X (Forward)};

  % Y axis (Right / East)
  \draw[axis, axisY] (0,0) -- ({\AL*\Yx},{\AL*\Yy});
  \node[font=\Large\bfseries, axisY, above right]
    at ({\AL*\Yx},{\AL*\Yy}) {Y (Right)};

  % Z axis (Down)
  \draw[axis, axisZ] (0,0) -- ({\AL*\Zx},{\AL*\Zy});
  \node[font=\Large\bfseries, axisZ, right]
    at ({\AL*\Zx},{\AL*\Zy}) {Z (Down)};

  % === Rotation arrows (elliptical arcs wrapping around each axis) ===

  % Roll - rotation around X axis
  % Arc lies in the Y-Z plane, centered on X axis
  % Projected semi-axes: Y direction (1,0) and Z direction (0,-1)
  \pgfmathsetmacro{\Rcx}{2.6*\Xx}  % center on X axis
  \pgfmathsetmacro{\Rcy}{2.6*\Xy}
  \pgfmathsetmacro{\Rr}{0.9}        % arc radius
  \draw[rotarrow, axisX!70!black]
    ({\Rcx+\Rr*\Yx*cos(200)+\Rr*\Zx*sin(200)},
     {\Rcy+\Rr*\Yy*cos(200)+\Rr*\Zy*sin(200)})
    \foreach \a in {200,210,...,490} {
      -- ({\Rcx+\Rr*\Yx*cos(\a)+\Rr*\Zx*sin(\a)},
          {\Rcy+\Rr*\Yy*cos(\a)+\Rr*\Zy*sin(\a)})
    };
  \node[font=\normalsize\bfseries, axisX!70!black, above]
    at ({\Rcx},{\Rcy+\Rr+0.15}) {Roll};

  % Pitch - rotation around Y axis
  % Arc lies in the X-Z plane, centered on Y axis
  % Projected semi-axes: X direction (-0.71,0.5) and Z direction (0,-1)
  %
  % Right-hand cyclic rule for a positive rotation about axis A: the arc
  % must be parametrised as (first-of-remaining-pair)*cos + (second)*sin
  % following the cyclic order X->Y->Z->X. Roll (about X) uses Y*cos+Z*sin
  % and yaw (about Z) uses X*cos+Y*sin, so pitch (about Y) must use
  % Z*cos+X*sin -- NOT X*cos+Z*sin, which would reverse the sense.
  % 右手系の巡回則: 軸Aまわりの正回転の円弧は、X->Y->Z->Xの巡回順で
  % 「（残り2軸のうち先の軸）*cos + （後の軸）*sin」としてパラメータ化する。
  % ロール（X軸まわり）はY*cos+Z*sin、ヨー（Z軸まわり）はX*cos+Y*sinを
  % 使うので、ピッチ（Y軸まわり）はZ*cos+X*sinでなければならない
  % （X*cos+Z*sinでは回転の向きが逆になる）。
  \pgfmathsetmacro{\Pcx}{2.6*\Yx}
  \pgfmathsetmacro{\Pcy}{2.6*\Yy}
  \pgfmathsetmacro{\Pr}{0.9}
  \draw[rotarrow, axisY!70!black]
    ({\Pcx+\Pr*\Zx*cos(200)+\Pr*\Xx*sin(200)},
     {\Pcy+\Pr*\Zy*cos(200)+\Pr*\Xy*sin(200)})
    \foreach \a in {200,210,...,490} {
      -- ({\Pcx+\Pr*\Zx*cos(\a)+\Pr*\Xx*sin(\a)},
          {\Pcy+\Pr*\Zy*cos(\a)+\Pr*\Xy*sin(\a)})
    };
  \node[font=\normalsize\bfseries, axisY!70!black, above right]
    at ({\Pcx+\Pr*0.3},{\Pcy+\Pr*0.7}) {Pitch};

  % Yaw - rotation around Z axis
  % Arc lies in the X-Y plane, centered on Z axis
  % Projected semi-axes: X direction (-0.71,0.5) and Y direction (1,0)
  \pgfmathsetmacro{\Wcx}{2.6*\Zx}
  \pgfmathsetmacro{\Wcy}{2.6*\Zy}
  \pgfmathsetmacro{\Wr}{0.9}
  \draw[rotarrow, axisZ!70!black]
    ({\Wcx+\Wr*\Xx*cos(160)+\Wr*\Yx*sin(160)},
     {\Wcy+\Wr*\Xy*cos(160)+\Wr*\Yy*sin(160)})
    \foreach \a in {160,170,...,450} {
      -- ({\Wcx+\Wr*\Xx*cos(\a)+\Wr*\Yx*sin(\a)},
          {\Wcy+\Wr*\Xy*cos(\a)+\Wr*\Yy*sin(\a)})
    };
  % Label pushed further outside the loop than the 0.8/0.3 multipliers
  % originally used (found by visual review to have its "w" overlapped
  % by the arc stroke).
  % ラベルを元の0.8/0.3倍率より外側へ押し出す（目視レビューで"w"の文字に
  % 弧の線が重なっていたのを修正）。
  \node[font=\normalsize\bfseries, axisZ!70!black, above left]
    at ({\Wcx-\Wr*1.4},{\Wcy+\Wr*0.55}) {Yaw};

\end{tikzpicture}
\end{document}
